\documentclass[12pt]{article}
\usepackage[utf8]{inputenc}
\usepackage[T1]{fontenc}
\usepackage{hyperref}
\title{Ta\'ino Language Resources for K-12 Curriculum Development}
\author{Compiled by Claude Sonnet 5 (claude-sonnet-5) for Andr\'as Kornai}
\date{}
\begin{document}
\maketitle

\section{Introduction}
Ta\'ino (ISO 639-3 tnq) is a dormant Arawakan language of the Caribbean. It has no native speakers. Sources disagree on its status: some call it extinct, while others describe reconstruction efforts and a small number of learners. The surviving record consists mostly of chronicler word lists, mainly nouns and loanwords.

\section{Primary Educational Resources}
No school curriculum or textbook was found in the search results. The UPR InDiLPE Caribe list mentions a Yale capstone titled ``Teaching Taino''. I did not open it.

\section{Linguistic and Reference Materials}
\begin{itemize}
\item UPR InDiLPE Caribe, list of works and projects on Ta\'ino (the group does not endorse the projects it lists): \url{https://generales.uprrp.edu/indilpe-caribe/?p=220}
\item The list as a PDF: \url{https://generales.uprrp.edu/indilpe-caribe/wp-content/uploads/sites/28/2024/03/List-of-works-and-projects-on-Taino.pdf}
\item Omniglot overview: \url{https://www.omniglot.com/writing/taino.htm}
\end{itemize}

\section{Religious and Community Resources}
\begin{itemize}
\item Casa Areyto, which says it works on language and culture revitalization, is mentioned on the IDIL resources page. I did not verify it further: \url{https://fr.idil2022-2032.org/resources/page/8/}
\item Liga Guak\'ia Ta\'ina-K\'e is mentioned in a student blog (advocacy material, to be verified): \url{https://hawksites.newpaltz.edu/decolonize/category/student-blog/page/2/}
\end{itemize}

\section{Digital and Audio Resources}
Tainona\'iki is a reconstructed variant by Javier A. Hern\'andez (2018), described on the IDIL page: \url{https://idil2022-2032.org/resources/page/5/}. Its linguistic basis is unverified.

\section{Curriculum Development Context}
A heritage curriculum can only be built on reconstruction, using the chronicler vocabulary and the revitalization projects above. Their quality varies. Any curriculum should state clearly that it is reconstructed. This dossier is short because little verified material exists.
\end{document}